% Esercitazione TEEA 3MAT — Argomento 3: resistività, temperatura, potenza ed energia
% Per rigenerare il PDF (dalla cartella strumenti, due volte):
%   pdflatex -output-directory=../../../materiali -jobname=esercitazione-resistivita-potenza-energia-3MAT esercitazione.tex
% ---------------------------------------------------------------------------
%  Stile comune di esercitazione e soluzioni (pdfLaTeX).
%  Viene incluso da esercitazione.tex e soluzioni.tex con \input.
% ---------------------------------------------------------------------------
\documentclass[11pt,a4paper]{article}
\usepackage[a4paper, top=14mm, bottom=18mm, left=16mm, right=16mm, footskip=9mm]{geometry}
\usepackage[utf8]{inputenc}
\usepackage[T1]{fontenc}
\usepackage{helvet}
\renewcommand{\familydefault}{\sfdefault}
\usepackage{amsmath, amssymb, graphicx, xcolor, fancyhdr, lastpage, tabularx, booktabs, array, enumitem}
\usepackage[most]{tcolorbox}
\usepackage[siunitx, american resistors, american currents]{circuitikz}
\ctikzset{bipoles/thickness=1, current/distance=0.35}
\usetikzlibrary{calc, arrows.meta}

\definecolor{titolo}{HTML}{2A7FA8}
\definecolor{grigio}{HTML}{E9ECEF}
\definecolor{tenue}{HTML}{555555}
\setlength{\parindent}{0pt}
\setlength{\parskip}{3pt}

\newcommand{\U}[1]{\,\mathrm{#1}}
\newcommand{\virg}{{,}}
\newcommand{\riga}[1][28mm]{\underline{\hspace{#1}}}

% intestazione come nelle esercitazioni della scuola
\newcommand{\intestazione}[2]{%
  \begin{minipage}[c]{20mm}\includegraphics[width=18mm]{logo-isis.png}\end{minipage}%
  \begin{minipage}[c]{\dimexpr\linewidth-20mm}
    {\fontsize{20}{24}\selectfont\bfseries\color{titolo} #1}\par\vspace{1mm}
    {\footnotesize\bfseries Indirizzo: IP14 -- MANUTENZIONE E ASSISTENZA TECNICA \quad -- \quad Classe 3MAT}
  \end{minipage}\par\vspace{2mm}
  {\color{titolo}\rule{\linewidth}{0.8pt}}\par\vspace{3mm}
  \begin{tcolorbox}[colback=grigio, colframe=grigio, arc=0pt, left=6pt, top=4pt, bottom=4pt]
    {\Large\bfseries #2}
  \end{tcolorbox}}

\newcommand{\nota}[1]{\begin{tcolorbox}[colback=grigio!45, colframe=grigio, arc=0pt, left=6pt, top=3pt, bottom=3pt]\small #1\end{tcolorbox}}
\newcommand{\tipo}[1]{\par\vspace{2mm}{\large\color{titolo} #1}\par\vspace{-1mm}{\color{titolo!60}\rule{\linewidth}{0.4pt}}\par\vspace{1mm}}
\newcommand{\esercizio}[2]{\par\vspace{3mm}\textbf{Esercizio #1.} \textbf{#2}\par}

% alimentazione: barra con etichetta
\newcommand{\rail}[2]{\draw[line width=1.2pt] ($#1+(-0.6,0)$) -- ($#1+(0.6,0)$); \node[above] at #1 {#2};}

\pagestyle{fancy}
\fancyhf{}
\renewcommand{\headrulewidth}{0pt}
\newcommand{\piede}[1]{\fancyfoot[L]{\scriptsize I.S.I.S. Sacile-Brugnera | sede: ``Piero Della Valentina'' | #1}
  \fancyfoot[R]{\scriptsize Pagina \thepage\ di \pageref{LastPage}}}

% Schemi degli esercizi (usati sia nell'esercitazione sia nelle soluzioni)
\definecolor{rame}{HTML}{C46A2B}

\newcommand{\lineaCarico}[3]{% #1 = V, #2 = R_L, #3 = R_C
\begin{circuitikz}[scale=0.8, transform shape, line width=0.6pt]
  \rail{(0,4.4)}{$V = #1$}
  \draw (0,4.4) -- (0,4.2) to[R, l_={$R_L = #2$}, i>^={$I$}] (0,2.2) to[short, *-] (0,2.2)
        to[R, l_={$R_C = #3$}] (0,0.2) -- (0,0) node[sground]{};
  \draw (0,2.2) to[short, -o] (1.4,2.2) node[right]{X};
\end{circuitikz}}

\newcommand{\partitoreI}[3]{% #1 = I, #2 = R1, #3 = R2
\begin{circuitikz}[scale=0.8, transform shape, line width=0.6pt]
  \draw (0,4.4) node[ocirc]{} to[short, i>^={$I = #1$}] (0,3.2);
  \draw (-1.6,3.2) -- (1.6,3.2); \draw (0,3.2) node[circ]{};
  \draw (-1.6,3.2) to[short, i>^={$I_1$}] (-1.6,2.6) to[R, l_={$R_1 = #2$}] (-1.6,0.6) -- (-1.6,0.3) -- (1.6,0.3);
  \draw (1.6,3.2) to[short, i>_={$I_2$}] (1.6,2.6) to[R, l^={$R_2 = #3$}] (1.6,0.6) -- (1.6,0.3);
  \draw (0,0.3) node[circ]{} -- (0,0.1) node[sground]{};
\end{circuitikz}}

\newcommand{\sezionePista}[2]{% #1 = w, #2 = t
\begin{tikzpicture}[line width=0.6pt]
  \fill[rame!70] (0,0) rectangle (3,0.6);
  \draw[rame!60!black, line width=0.9pt] (0,0) rectangle (3,0.6);
  \draw[{Latex[length=2mm]}-{Latex[length=2mm]}] (0,-0.3) -- node[below]{$w = #1$} (3,-0.3);
  \draw[{Latex[length=2mm]}-{Latex[length=2mm]}] (3.3,0) -- node[right]{$t = #2$} (3.3,0.6);
\end{tikzpicture}}

\piede{Esercitazione TEEA 3MAT}

\begin{document}

\intestazione{ESERCITAZIONE DI TEEA}{Resistività, temperatura, potenza ed energia}

\nota{Esercizi di allenamento su resistenza di un conduttore, rame e alluminio, piste dei circuiti stampati,
resistenza e temperatura, effetto Joule, computo dei consumi e partitori. Scrivi sempre la formula, sostituisci
i valori con le unità di misura e controlla il risultato. Dati: $\rho_{Cu} = 0{,}0175\U{\Omega \cdot mm^2/m}$,
$\rho_{Al} = 0{,}028\U{\Omega \cdot mm^2/m}$, $\alpha_{Cu} = 0{,}0039\U{1/^\circ C}$, densità rame $8{,}96\U{g/cm^3}$,
alluminio $2{,}70\U{g/cm^3}$. L'energia va espressa in Wh o kWh.}

\vspace{2mm}
\textbf{COGNOME E NOME:} \riga[95mm] \hfill \textbf{CLASSE:} \riga[22mm]

\tipo{Resistenza di un conduttore}

\esercizio{1}{Un cavo di rame è lungo $50\U{m}$ e ha sezione $1{,}5\U{mm^2}$.
a) Calcola la sua resistenza. b) Quale sezione deve avere un cavo di alluminio della stessa lunghezza per avere
la stessa resistenza? c) Calcola il diametro dei due conduttori. d) Quanto pesano i due cavi?}
\vspace{2mm}
\renewcommand{\arraystretch}{2.1}
\begin{tabular}{@{}l l@{}}
a) $R_{Cu} =$ \riga & b) $S_{Al} =$ \riga\\
c) $d_{Cu} =$ \riga & \phantom{b)} $d_{Al} =$ \riga\\
d) peso rame $=$ \riga & \phantom{b)} peso alluminio $=$ \riga\\
\end{tabular}
\vspace{22mm}

\esercizio{2}{Un conduttore di una linea aerea in alluminio è lungo $2\U{km}$ e ha sezione $150\U{mm^2}$.
a) Calcola la sua resistenza. b) Quale sezione di rame darebbe la stessa resistenza?
c) Calcola quanto pesa $1\U{km}$ di conduttore in alluminio e $1\U{km}$ di conduttore in rame.
d) Spiega in due righe perché le linee aeree si fanno in alluminio.}
\vspace{2mm}
\begin{tabular}{@{}l l@{}}
a) $R_{Al} =$ \riga & b) $S_{Cu} =$ \riga\\
c) $1\U{km}$ di alluminio $=$ \riga & \phantom{b)} $1\U{km}$ di rame $=$ \riga\\
\end{tabular}\par
d) \riga[165mm]\par\vspace{3mm}\riga[172mm]
\vspace{14mm}

\newpage
\esercizio{3}{Una pista di rame di un circuito stampato è lunga $80\U{mm}$, larga $0{,}8\U{mm}$ e spessa $35\U{\mu m}$.
a) Calcola la sezione e la resistenza della pista. b) Nella pista passano $3\U{A}$: calcola la tensione persa
ai suoi capi e la potenza dissipata. c) Quanto dovrebbe essere larga la pista per avere metà resistenza?}
\begin{minipage}[t]{0.36\linewidth}\vspace{4mm}\centering
\sezionePista{0{,}8\,\mathrm{mm}}{35\,\mu\mathrm{m}}
\end{minipage}\hfill
\begin{minipage}[t]{0.60\linewidth}\vspace{2mm}
\renewcommand{\arraystretch}{2.1}
\begin{tabular}{@{}l l@{}}
a) $S =$ \riga & $R =$ \riga\\
b) $V =$ \riga & $P =$ \riga\\
c) $w =$ \riga & \\
\end{tabular}
\end{minipage}
\vspace{16mm}

\tipo{Resistenza e temperatura, effetto Joule}

\esercizio{4}{L'avvolgimento in rame di un motore ha una resistenza di $4\U{\Omega}$ a $20\U{^\circ C}$.
a) Quanto vale la resistenza a $90\U{^\circ C}$? b) Dopo un'ora di lavoro il manutentore misura con
l'ohmmetro una resistenza di $4{,}8\U{\Omega}$: a che temperatura è arrivato l'avvolgimento?}
\vspace{2mm}
\renewcommand{\arraystretch}{2.1}
\begin{tabular}{@{}l l@{}}
a) $R_{90} =$ \riga & \\
b) $\Delta T =$ \riga & $T =$ \riga\\
\end{tabular}
\vspace{24mm}

\esercizio{5}{In un cavo di rame con resistenza $0{,}5\U{\Omega}$ a $20\U{^\circ C}$ passa una corrente di $16\U{A}$.
a) Calcola la potenza dissipata nel cavo a $20\U{^\circ C}$ e il calore prodotto in $1\U{h}$, prima in joule e poi in Wh.
b) Il cavo si scalda fino a $60\U{^\circ C}$: calcola la nuova resistenza e la nuova potenza.
c) Calcola l'energia persa in $8\U{h}$ nei due casi. La potenza è aumentata o diminuita? Perché?}
\vspace{2mm}
\begin{tabular}{@{}l l l@{}}
a) $P_{20} =$ \riga & $Q =$ \riga[24mm] J & $=$ \riga[20mm] Wh\\
b) $R_{60} =$ \riga & $P_{60} =$ \riga & \\
c) $E_{20} =$ \riga & $E_{60} =$ \riga & \\
\end{tabular}\par
\riga[172mm]
\vspace{8mm}

\newpage
\tipo{Computo energetico e spesa}

\esercizio{6}{Completa la tabella dei consumi di un appartamento (un mese = $30$ giorni) e calcola la spesa
mensile per l'energia con un prezzo di $0{,}25$ euro/kWh. Quale apparecchio consuma di più?}
\vspace{2mm}
{\renewcommand{\arraystretch}{1.9}
\begin{tabularx}{\linewidth}{@{}l c c >{\centering\arraybackslash}X >{\centering\arraybackslash}X@{}}
\toprule
\textbf{Apparecchio} & \textbf{Potenza} & \textbf{Uso al giorno} & \textbf{Energia al giorno} & \textbf{Energia al mese}\\
\midrule
bollitore & $2000\U{W}$ & $15\U{min}$ & & \\
scaldabagno & $1200\U{W}$ & $2\U{h}$ & & \\
4 lampade LED & $8\U{W}$ ciascuna & $6\U{h}$ & & \\
computer & $200\U{W}$ & $4\U{h}$ & & \\
\midrule
\textbf{totale} & & & & \\
\bottomrule
\end{tabularx}}
\vspace{2mm}

spesa mensile $=$ \riga[30mm] euro \qquad consuma di più: \riga[45mm]
\vspace{6mm}

\tipo{Partitori e consumi}

\esercizio{7}{Un impianto in continua a $24\U{V}$ alimenta un carico con $R_C = 4{,}45\U{\Omega}$ attraverso un cavo
di rame da $1{,}5\U{mm^2}$ con $15\U{m}$ di andata e $15\U{m}$ di ritorno.
a) Calcola la resistenza della linea $R_L$. b) Con il partitore di tensione calcola la tensione sul carico e quella
persa sulla linea. c) Calcola la corrente, le potenze e il rendimento. d) Calcola l'energia data al carico e quella
persa nella linea in $6\U{h}$. e) Ripeti b) e il rendimento con un cavo da $2{,}5\U{mm^2}$.}
\begin{minipage}[t]{0.34\linewidth}\vspace{0pt}\centering
\lineaCarico{24\,\mathrm{V}}{\,?}{4{,}45\,\Omega}
\end{minipage}\hfill
\begin{minipage}[t]{0.64\linewidth}\vspace{2mm}
\renewcommand{\arraystretch}{2.0}
\begin{tabular}{@{}l l@{}}
a) $R_L =$ \riga & \\
b) $V_C =$ \riga & $V_L =$ \riga\\
c) $I =$ \riga & $\eta =$ \riga\\
\phantom{c)} $P_C =$ \riga & $P_L =$ \riga\\
d) $E_C =$ \riga & $E_L =$ \riga\\
e) $V_C =$ \riga & $\eta =$ \riga\\
\end{tabular}
\end{minipage}

\newpage
\esercizio{8}{Due carichi resistivi in parallelo, $R_1 = 8\U{\Omega}$ e $R_2 = 24\U{\Omega}$, sono attraversati in tutto da
$I = 6\U{A}$ e funzionano $4\U{h}$ al giorno per $30$ giorni.
a) Con il partitore di corrente calcola $I_1$ e $I_2$ e la tensione ai capi. b) Calcola la potenza di ciascun carico.
c) Calcola l'energia consumata nel mese da ciascun carico, in kWh, e il costo totale con $0{,}25$ euro/kWh.
d) Quale dei due carichi consuma di più? Perché?}
\begin{minipage}[t]{0.42\linewidth}\vspace{0pt}\centering
\partitoreI{6\,\mathrm{A}}{8\,\Omega}{24\,\Omega}
\end{minipage}\hfill
\begin{minipage}[t]{0.56\linewidth}\vspace{2mm}
\renewcommand{\arraystretch}{2.1}
\begin{tabular}{@{}l l@{}}
a) $I_1 =$ \riga & $I_2 =$ \riga\\
\phantom{a)} $V =$ \riga & \\
b) $P_1 =$ \riga & $P_2 =$ \riga\\
c) $E_1 =$ \riga & $E_2 =$ \riga\\
\phantom{c)} costo $=$ \riga & \\
\end{tabular}
\end{minipage}\par
d) \riga[165mm]\par\vspace{3mm}\riga[172mm]

\end{document}
